\section{Item Review}

Public examples are treated as source leads, not as ground truth. A post,
article, repository, or paper can motivate an archetype, but the benchmark item
must be represented independently, checked for ambiguity, and tied to an
objective answer before it can enter a paper split.

\paragraph{Item-card lifecycle.}
Each item card records the item ID, source summary, answer derivation, expected
answer, scorer contract, ambiguity notes, split policy, leakage risk, review
status, reviewer, and review date. Cards with incomplete review text cannot
support paper claims.

\paragraph{Split policy.}
The repo separates public examples, development data, paper-frozen data,
private refreshes, and canary items. Public examples are useful for
transparency and failure galleries, but public exposure makes them weak
evidence for contamination-resistant ranking.

\paragraph{Human-validation policy.}
For the fast preprint, human-triviality is a design target rather than a
measured result. The paper therefore relies on reviewed item cards, answer
derivations, ambiguity notes, source-policy checks, and scorer-gold fixtures.
It does not report human accuracy or response-time distributions. A measured
human baseline remains future validation after the split is frozen.
